\section{Related Work}
\paragraph{Steering and refusal.} Learned interventions on frozen-model activations include representation finetuning and bidirectional preference steering \citep{wu2024reft,cao2024bipo}. We audit one fixed-vector injection protocol. Related work separates steering efficacy from the preservation of other behaviors under distribution shift \citep{goyal2026specificity}. \citet{arditi2024refusal} study a refusal direction in chat models. This supports evaluating harmful compliance and benign task utility together. Our learned vector, its residual decomposition, and its final-prompt-token injection protocol are specific to the setup below; we do not equate them with a universally shared refusal direction.

\paragraph{Random steering.} \citet{korznikov2026rogue} report that random fixed-vector steering can increase harmful compliance from zero to 1--13\% in the revised version. Their intervention persists across prompt and generated tokens, whereas our intervention is a single injection at the final prompt token. Our attack prompts are also selected for prior success. Random perturbations move our judge score in the opposite direction. The protocol and input-selection differences are plausible explanations, not a tested mechanism or a direct numerical replication.

\paragraph{Attack and benign evaluations.} PAIR iteratively constructs semantic jailbreaks through black-box interaction \citep{chao2023pair}. HarmBench \citep{mazeika2024harmbench} and JailbreakBench \citep{chao2024jbb} emphasize standardized attack evaluation. WildTeaming introduces WildJailbreak with both vanilla and adversarial harmful and benign examples \citep{jiang2024wildteaming}. Its benign requests that resemble harmful ones motivate our paired wrapper analysis. Because our benign intervention and evaluation share this source, that intervention tests in-distribution data composition rather than generalization to arbitrary wrapper styles.

\paragraph{Scope of the audit.} Prior work compares automated jailbreak evaluators and studies judge calibration and robustness \citep{gao2026judge,mu2026evaluators}. We complement these evaluations by distinguishing a harmfulness judge flip from an observed refusal, and an observed refusal from a calibrated utility failure. We use matched baselines and explicit reading scopes to expose cases where those distinctions change reported counts. Numerical agreement across implementations is an important reproducibility check, but it does not validate the judge on a new distribution or language.
